% Rapport de l'étude « désordre et itérations de Lloyd » — produit par ./run.sh (latexmk -lualatex rapport.tex)
\documentclass[10pt,a4paper]{article}

\usepackage{fontspec}
\usepackage[french]{babel}
\usepackage{amsmath,amssymb}
\usepackage{booktabs}
\usepackage{graphicx}
\usepackage{subcaption}
\usepackage{float}
\usepackage{pgfplots}
\pgfplotsset{compat=1.18}
\usepackage[margin=2.2cm]{geometry}
\usepackage{xcolor}
\usepackage{listings}
\usepackage{enumitem}
\usepackage[colorlinks=true,linkcolor=blue!50!black,urlcolor=blue!50!black]{hyperref}

\setlist{itemsep=1pt,topsep=3pt}
\emergencystretch=3em
\lstset{basicstyle=\ttfamily\footnotesize,columns=fullflexible,frame=single,framerule=0.3pt,
  backgroundcolor=\color{gray!6},breaklines=true,keepspaces=true,upquote=true}
\newcommand{\code}[1]{\texttt{#1}}
\newcommand{\vect}[1]{\boldsymbol{#1}}
\definecolor{cU}{HTML}{2D7DD2}
\definecolor{cR}{HTML}{E0524D}
\definecolor{cH}{HTML}{2E9E5B}

% courbes d'un indicateur pour les trois états initiaux (abscisse : itérations + 1)
\pgfplotsset{
  etude/.style={width=0.48\linewidth,height=4.6cm,xlabel={itérations $+1$},
    grid=major,grid style={gray!20},legend style={font=\scriptsize},
    every axis plot/.append style={thick,mark size=1.1pt,error bars/y dir=both,error bars/y explicit}},
}
\newcommand{\courbes}[2]{% #1 colonne, #2 options supplémentaires
  \begin{axis}[etude,xmode=log,#2]
    \addplot[cU,mark=*] table[x expr=\thisrow{iteration}+1,y=#1,y error=#1_sd]{data/uniform.dat};
    \addplot[cR,mark=square*] table[x expr=\thisrow{iteration}+1,y=#1,y error=#1_sd]{data/random.dat};
    \addplot[cH,mark=triangle*] table[x expr=\thisrow{iteration}+1,y=#1,y error=#1_sd]{data/hex.dat};
  \end{axis}}

\title{Désordre d'un pavage de Voronoï au fil des itérations de Lloyd\\[3pt]
  \large Petite étude réalisée avec \code{cellyod}}
\author{Dépôt \code{cellYOd}, répertoire \code{etudes/lloyd\_desordre}}
\date{\today}

\begin{document}
\maketitle

\begin{abstract}
  Les itérations de Lloyd rendent un pavage de Voronoï de plus en plus régulier, mais à quelle vitesse, et
  jusqu'où ? On suit ici, sans bundles ($\alpha=0$), neuf indicateurs de désordre (dispersion des aires et des
  longueurs de paroi, nombre de côtés, ordre orientationnel, énergie de Lloyd\dots) pour trois états initiaux, de
  0 à 1000 itérations, sur des tissus d'environ 2000 cellules et cinq graines. Le désordre décroît d'abord très
  vite (une dizaine d'itérations divise par trois la dispersion des aires d'un tirage uniforme), puis lentement,
  en loi de puissance de faible exposant : partant d'un tirage aléatoire, Lloyd ne retrouve pas le réseau
  hexagonal mais un état centroïdal avec des défauts topologiques (paires pentagone--heptagone) qui
  disparaissent très lentement. Les indicateurs se répartissent en deux familles, métrique et topologique, qui
  évoluent à des vitesses différentes ; deux nombres suffisent à décrire le désordre.
\end{abstract}

%======================================================================
\section{Question et dispositif}
%======================================================================

\textbf{Question.} À partir d'un ensemble de germes plus ou moins désordonné, l'itération de Lloyd déplace chaque
germe au centroïde de sa cellule de Voronoï. On veut caractériser le désordre du pavage en fonction du nombre
d'itérations $k$, choisir des indicateurs pertinents et savoir si l'état initial est oublié.

\textbf{Simulations.} \code{cellyod}, domaine carré $40\times40$ (longueurs en diamètres de cellule, environ
2000 cellules), sans bundles ni forme ($\alpha=0$ : Lloyd pur). Trois états initiaux :
\begin{itemize}
  \item \code{uniform} : tirage uniforme (processus de Poisson), le plus désordonné ;
  \item \code{random} : tirage avec une distance minimale de $0.7$ espacement (processus à cœur dur) ;
  \item \code{hex} : réseau hexagonal dont chaque germe est déplacé au hasard d'au plus $0.5$ pas
    (\code{disorder 0.5}), c'est-à-dire un cristal fortement perturbé.
\end{itemize}
Pour chaque état, 5 graines ; chaque calcul est enregistré aux points de contrôle
$k={}$0, 1, 2, 3, 5, 7, 10, 15, 20, 30, 50, 70, 100, 150, 200, 300, 500, 700 et 1000, en reprenant le fichier
précédent avec \code{-load} (ce qui donne exactement les mêmes germes qu'un calcul d'une traite). Les courbes
donnent la moyenne sur les 5 graines, les barres l'écart type.

\textbf{Reproduire l'étude.} Tout est produit par un script (environ une minute) :
\begin{lstlisting}
cd etudes/lloyd_desordre
./run.sh        # compile cellyod et mesure, simulations, mesures, figures, ce rapport
\end{lstlisting}
Les réglages (taille, états, graines, points de contrôle) sont en tête de \code{run.sh}. Les indicateurs sont
calculés par \code{mesure} (\code{mesure.cpp}, 140 lignes), qui relit les fichiers de cellules avec
le cœur de \code{cellyod} : \code{./mesure work/*.txt} donne une ligne d'indicateurs par fichier.

%======================================================================
\section{Indicateurs de désordre}
%======================================================================

Pour ne pas mêler désordre et effet du bord, seules les cellules \emph{profondes} sont comptées : ni cellule de
bord, ni voisine d'une cellule de bord (environ 1700 cellules). Sur un pavage hexagonal parfait, tous les
indicateurs prennent alors exactement leur valeur de référence (dernière ligne du tableau~\ref{tab:valeurs}).
\begin{itemize}
  \item \textbf{Métriques} :
    $\mathrm{cv}_A=\sigma_A/\langle A\rangle$, coefficient de variation des aires ;
    $\mathrm{cv}_L$, coefficient de variation des longueurs de paroi (parois entre deux cellules profondes) ;
    $\mathrm{cv}_{NN}$, coefficient de variation de la distance d'un germe à son plus proche voisin ;
    $\langle SI\rangle=\langle4\pi A/P^2\rangle$, indice de forme (hexagone : 0.9069) ;
  \item \textbf{Énergie} : moment d'inertie normalisé
    $G=\sum_i\int_{V_i}\|\vect x-\vect p_i\|^2d\vect x\,/\,(2n\langle A\rangle^2)$, l'énergie que minimise Lloyd,
    rendue sans dimension ; minimale pour le pavage hexagonal, $G_{\mathrm{hex}}=5/(36\sqrt3)=0.080188$ ;
  \item \textbf{Topologiques} : $\sigma_n$, écart type du nombre de côtés (0 pour un pavage hexagonal) ; $p_6$,
    fraction d'hexagones ; $\psi_6=\langle|\frac1{n_i}\sum_j e^{6i\theta_{ij}}|\rangle$, ordre orientationnel
    local, $\theta_{ij}$ étant l'angle du germe $j$, voisin de Voronoï de $i$ (1 pour un réseau hexagonal) ;
  \item \textbf{Convergence} : résidu $\langle\|\vect p_i-\vect c_i\|\rangle/s$, distance moyenne entre un germe et
    le centroïde de sa cellule ($s$ : espacement visé) ; il est nul au point fixe de Lloyd (tessellation
    centroïdale), que celle-ci soit ordonnée ou non.
\end{itemize}

\textbf{Contrôle.} Pour le tirage uniforme sans itération (pavage de Voronoï--Poisson), on mesure
$\mathrm{cv}_A=0.538$, $\sigma_n=1.341$ et $p_6=0.297$, proches des valeurs connues pour un domaine infini
(environ 0.53, 1.33 et 0.29).

%======================================================================
\section{Résultats}
%======================================================================

\begin{figure}[H]
  \centering
  \begin{subfigure}{0.24\linewidth}\includegraphics[width=\linewidth]{figures/vue_0.png}\caption{$k=0$}\end{subfigure}\hfill
  \begin{subfigure}{0.24\linewidth}\includegraphics[width=\linewidth]{figures/vue_3.png}\caption{$k=3$}\end{subfigure}\hfill
  \begin{subfigure}{0.24\linewidth}\includegraphics[width=\linewidth]{figures/vue_20.png}\caption{$k=20$}\end{subfigure}\hfill
  \begin{subfigure}{0.24\linewidth}\includegraphics[width=\linewidth]{figures/vue_200.png}\caption{$k=200$}\end{subfigure}
  \caption{Tirage \code{uniform} ($16\times16$) au fil des itérations, cellules colorées par nombre de côtés :
    \textcolor[HTML]{D64541}{4}, \textcolor[HTML]{F0963C}{5}, \textcolor[HTML]{8DB87A}{6},
    \textcolor[HTML]{5FAAE1}{7}, \textcolor[HTML]{2D6EB4}{8}, \textcolor[HTML]{193C6E}{9 et plus}. Les cellules du
    bord, coupées par le carré, ont naturellement moins de côtés. Tracé à partir de l'export polygones de
    \code{cellyod}.}
  \label{fig:vues}
\end{figure}

\begin{figure}[H]
  \centering
  \begin{tikzpicture}\courbes{cvA}{ymode=log,ylabel={$\mathrm{cv}_A$},
      legend entries={\code{uniform},\code{random},\code{hex}},legend pos=south west}\end{tikzpicture}\hfill
  \begin{tikzpicture}\courbes{cvL}{ymode=log,ylabel={$\mathrm{cv}_L$}}\end{tikzpicture}\\[1mm]
  \begin{tikzpicture}\courbes{sigmaN}{ymode=log,ylabel={$\sigma_n$}}\end{tikzpicture}\hfill
  \begin{tikzpicture}
    \begin{axis}[etude,xmode=log,ymode=log,ylabel={$1-p_6$}]
      \addplot[cU,mark=*] table[x expr=\thisrow{iteration}+1,y expr=1-\thisrow{p6},y error=p6_sd]{data/uniform.dat};
      \addplot[cR,mark=square*] table[x expr=\thisrow{iteration}+1,y expr=1-\thisrow{p6},y error=p6_sd]{data/random.dat};
      \addplot[cH,mark=triangle*] table[x expr=\thisrow{iteration}+1,y expr=1-\thisrow{p6},y error=p6_sd]{data/hex.dat};
    \end{axis}
  \end{tikzpicture}\\[1mm]
  \begin{tikzpicture}
    \begin{axis}[etude,xmode=log,ymode=log,ylabel={$G-G_{\mathrm{hex}}$}]
      \addplot[cU,mark=*] table[x expr=\thisrow{iteration}+1,y expr=\thisrow{G}-0.0801875,y error=G_sd]{data/uniform.dat};
      \addplot[cR,mark=square*] table[x expr=\thisrow{iteration}+1,y expr=\thisrow{G}-0.0801875,y error=G_sd]{data/random.dat};
      \addplot[cH,mark=triangle*] table[x expr=\thisrow{iteration}+1,y expr=\thisrow{G}-0.0801875,y error=G_sd]{data/hex.dat};
    \end{axis}
  \end{tikzpicture}\hfill
  \begin{tikzpicture}\courbes{residu}{ymode=log,ylabel={résidu germe--centroïde}}\end{tikzpicture}
  \caption{Indicateurs en fonction du nombre d'itérations (échelles logarithmiques ; moyenne et écart type sur 5
    graines, environ 1700 cellules profondes chacune). Bleu : \code{uniform} ; rouge : \code{random} ; vert :
    \code{hex} (\code{disorder 0.5}).}
  \label{fig:courbes}
\end{figure}

\begin{table}[H]
  \centering
  \input{data/tableau.tex}
  \caption{Valeurs moyennes (5 graines) à quelques points de contrôle.}
  \label{tab:valeurs}
\end{table}

\subsection{Deux régimes}

\textbf{Régime rapide} (les dix premières itérations). L'essentiel du désordre métrique disparaît : pour le
tirage uniforme, $\mathrm{cv}_A$ passe de 0.54 à 0.17, $\mathrm{cv}_{NN}$ de 0.53 à 0.11 et $G$ de 0.161 à 0.085,
soit plus de 90\,\% du chemin vers $G_{\mathrm{hex}}$. Les très petites cellules du tirage de Poisson, entourées de
grandes cellules, s'élargissent dès les premières itérations (figure~\ref{fig:vues}, $k=3$).

\textbf{Régime lent} (au-delà d'une vingtaine d'itérations). Les indicateurs continuent de décroître, mais en loi
de puissance de faible exposant. Entre $k=20$ et $k=1000$, pour le tirage uniforme,
$\mathrm{cv}_A\propto k^{-0.27}$, $\mathrm{cv}_L\propto k^{-0.18}$, $\sigma_n\propto k^{-0.14}$,
$1-p_6\propto k^{-0.28}$ et $G-G_{\mathrm{hex}}\propto k^{-0.38}$ (ajustements sur les moyennes). Après 1000
itérations, il reste 14\,\% de cellules non hexagonales. Dans le même temps, le résidu germe--centroïde décroît
en $k^{-0.96}$, soit comme $1/k$ : les germes sont de plus en plus proches d'une tessellation centroïdale, mais
celle-ci contient des défauts.

\textbf{Nature des défauts.} Les cellules non hexagonales restantes sont uniquement des pentagones et des
heptagones, en nombres égaux (pour la graine 1 du tirage uniforme : 9.7\,\% de chaque à $k=200$, 7.9\,\% à
$k=1000$), le plus souvent accolés en paires (figure~\ref{fig:vues}, $k=200$) : ce sont des dislocations du
réseau hexagonal, et leurs
alignements forment des joints entre domaines hexagonaux d'orientations différentes ($\psi_6=0.84$ à $k=1000$).
Une dislocation est une configuration quasi stationnaire pour Lloyd (le résidu est faible) : elle ne disparaît
qu'en se déplaçant jusqu'au bord ou en s'annihilant avec une dislocation opposée, d'où la lenteur du second
régime. Lloyd se comporte comme une descente d'énergie qui se bloque dans des états métastables, à la
manière d'une trempe.

\subsection{Mémoire de l'état initial}

Les tirages \code{uniform} et \code{random} partent de niveaux de désordre très différents ($\mathrm{cv}_A$ 0.54
et 0.19) mais leurs courbes se rejoignent : à $k=100$, $p_6$ vaut 0.77 et 0.78, $\sigma_n$ 0.48 et 0.47 ; à
$k=1000$, $\mathrm{cv}_A$ vaut 0.046 et 0.044. La topologie oublie l'état initial en une centaine d'itérations ; la
dispersion des aires en garde un peu plus longtemps la trace.

Le cristal perturbé (\code{hex}, \code{disorder 0.5}) part pourtant d'un désordre topologique comparable à
\code{random} ($\sigma_n=0.92$, $p_6=0.44$) mais suit une tout autre trajectoire : les germes « se souviennent » du
réseau sous-jacent, et le pavage redevient presque parfaitement hexagonal en une cinquantaine d'itérations
($p_6=0.97$, $\psi_6=0.96$). Il reste alors sur un plateau : quelques dislocations subsistent, vraisemblablement
parce que le réseau hexagonal ne s'accorde pas au carré $40\times40$ : ses rangées sont espacées de $0.825$,
soit 48.5 rangées sur la hauteur, un nombre non entier. Entre $k=150$ et $k=1000$, $\sigma_n$ remonte légèrement (de 0.152 à 0.166) pour les cinq graines,
un effet de l'ordre de l'écart type entre graines ; une explication possible, non vérifiée ici, est la migration
de dislocations du bord vers les cellules profondes.

\subsection{Quels indicateurs retenir ?}

\begin{figure}[H]
  \centering
  \begin{tikzpicture}
    \begin{axis}[width=0.46\linewidth,height=5cm,xlabel={$\sigma_n$},ylabel={$\mathrm{cv}_L$},grid=major,
        grid style={gray!20},only marks,mark size=1.3pt]
      \addplot[cU,mark=*] table[x=sigmaN,y=cvL]{data/uniform.dat};
      \addplot[cR,mark=square*] table[x=sigmaN,y=cvL]{data/random.dat};
      \addplot[cH,mark=triangle*] table[x=sigmaN,y=cvL]{data/hex.dat};
    \end{axis}
  \end{tikzpicture}\hfill
  \begin{tikzpicture}
    \begin{axis}[width=0.46\linewidth,height=5cm,xlabel={$\sigma_n$},ylabel={$\mathrm{cv}_A$},grid=major,
        grid style={gray!20},only marks,mark size=1.3pt,ymode=log]
      \addplot[cU,mark=*] table[x=sigmaN,y=cvA]{data/uniform.dat};
      \addplot[cR,mark=square*] table[x=sigmaN,y=cvA]{data/random.dat};
      \addplot[cH,mark=triangle*] table[x=sigmaN,y=cvA]{data/hex.dat};
    \end{axis}
  \end{tikzpicture}
  \caption{Indicateurs tracés l'un en fonction de l'autre, pour tous les états et tous les points de contrôle.
    À gauche, $\mathrm{cv}_L$ et $\sigma_n$ décrivent la même chose, quel que soit l'état initial ; à droite,
    $\mathrm{cv}_A$ et $\sigma_n$ ne sont pas liés de la même façon selon l'état initial.}
  \label{fig:correlations}
\end{figure}

Sur l'ensemble des points (3 états, 19 points de contrôle), les indicateurs forment deux familles
(figure~\ref{fig:correlations}, coefficients de corrélation linéaire $r$) :
\begin{itemize}
  \item \textbf{topologique} : $\sigma_n$, $p_6$, $\psi_6$ et, plus surprenant, $\mathrm{cv}_L$ sont presque
    équivalents ($|r|\ge0.96$ deux à deux, $0.985$ entre $\mathrm{cv}_L$ et $\sigma_n$) : la dispersion des
    longueurs de paroi vient vraisemblablement surtout des parois courtes qui entourent les défauts. Ils mesurent la densité de défauts et évoluent lentement ;
  \item \textbf{métrique} : $\mathrm{cv}_A$, $\mathrm{cv}_{NN}$, $G$ et $\langle SI\rangle$ mesurent l'hétérogénéité
    des tailles et des formes ($|r|\ge0.94$ entre $\mathrm{cv}_A$, $\mathrm{cv}_{NN}$ et $\langle SI\rangle$) ; ils
    chutent pendant le régime rapide et sont moins liés à la topologie ($r=0.89$ entre $\mathrm{cv}_A$ et
    $\sigma_n$, $0.72$ entre $G$ et $\sigma_n$, avec une relation qui dépend de l'état initial).
\end{itemize}
Pour décrire le désordre d'un pavage, un couple suffit : \textbf{un indicateur métrique}, de préférence
$\mathrm{cv}_A$ (le plus simple, déjà calculé par \code{cellyod}) ou $G-G_{\mathrm{hex}}$ (lié directement à
l'énergie de Lloyd), et \textbf{un indicateur topologique}, $\sigma_n$ ou $1-p_6$ (simples, sans paramètre).
$\psi_6$ apporte en plus l'information d'orientation (domaines et joints), mais demande les voisinages.
Le résidu germe--centroïde n'est pas un indicateur de désordre : il mesure la distance au point fixe.

%======================================================================
\section{Conclusions et pistes}
%======================================================================

\begin{itemize}
  \item Lloyd régularise très vite les tailles et les formes (une dizaine d'itérations), puis élimine lentement
    les défauts topologiques, en loi de puissance de faible exposant ; à partir d'un tirage aléatoire, il ne
    retrouve pas le réseau hexagonal, même après 1000 itérations ($p_6\approx0.86$).
  \item Le nombre d'itérations est donc un réglage du désordre : quelques itérations pour un tissu encore varié,
    quelques dizaines pour des cellules régulières mais une topologie désordonnée ; pour un tissu presque
    hexagonal, mieux vaut partir d'un réseau perturbé (\code{hex}) que d'un tirage aléatoire.
  \item Deux nombres suffisent à décrire le désordre : un métrique ($\mathrm{cv}_A$ ou $G$) et un topologique
    ($\sigma_n$ ou $p_6$).
\end{itemize}
\textbf{Pistes.} Effet de la taille du domaine (le plateau de \code{hex} dépend vraisemblablement de
l'incompatibilité entre le réseau et le carré ; à tester avec plusieurs tailles, ou avec une forme circulaire
via \code{shape}) ; plus de graines pour préciser les exposants ; recuit (petite perturbation des germes puis
Lloyd) pour éliminer les dislocations ; mêmes mesures avec des bundles, pour séparer le désordre imposé par
l'attraction de celui hérité de l'état initial.

\end{document}
